% Section 1 — Introduction
\section{Introduction}
\label{sec:intro}

\paragraph{The raw network} In 2016 AlphaGo won with a move that its own imitation network rated at
about one chance in ten thousand; the move came from search, paid for at the moment of play
\citep{silver2016alphago,hassabis2016seoul}. A year later AlphaGo Zero learned Go from self-play
alone, and its paper contains a quieter result that matters more here. Playing the network's
preferred move with no search at all, AlphaGo Zero rated about 3,000 Elo---close to the version that
had first beaten a professional---against 5,185 with search \citep{silver2017zero}. Much of what search
had discovered across millions of games had been \emph{amortized} into a single proposal. AlphaZero
and MuZero carried the recipe to other games and to learned rules
\citep{silver2018alphazero,schrittwieser2020muzero}.

\paragraph{Innovation still pays at the moment of the problem} Machine discovery has since moved from
games to science and engineering, and it has kept paying for search per problem. AlphaTensor found a
47-multiplication scheme for $4\times4$ matrices in arithmetic modulo 2 \citep{fawzi2022alphatensor};
AlphaDev found sorting routines merged into the LLVM C\texttt{++} library \citep{mankowitz2023alphadev};
FunSearch improved bounds on cap sets by pairing a language model with an evaluator
\citep{romeraparedes2024funsearch}; AlphaEvolve extended that loop to matrix multiplication,
scheduling and, by 2026, to cache policies, circuit design and commercial optimization
\citep{novikov2025alphaevolve,deepmind2026alphaevolve}. Automated research systems close the loop
around whole experiments \citep{lu2024aiscientist,recursive2026flint}, and Socher's \emph{Eureka
Machine} describes swarms of agents proposing and testing ideas across all of science
\citep{socher2026eureka}. Engineering and laboratory practice look the same: a bracket is found by
topology optimization over hundreds of simulations \citep{sigmund2001topopt}, and a reaction
condition by Bayesian optimization over a campaign of experiments
\citep{snoek2012bo,shields2021reaction}. Each system is powerful. Each starts every new problem from
its search, not from what earlier searches learned.

\paragraph{Innovation as a game} This paper begins from an observation and a hypothesis. The
observation is that many innovation problems already have the anatomy of a game
(\cref{fig:game}). There is a \emph{state}: the requirements, observations and budget that define the
problem. There are \emph{moves}: complete candidate solutions written in a typed language---a learning
procedure, a next experiment, a CAD construction program. There are \emph{rules}: the language's
guarantees and the problem's hard constraints. There is an \emph{outcome}, measured by an independent
verifier rather than by anyone's opinion. And there is an \emph{opponent} worth beating: the best
verified solution known for that problem, which rises every time it is beaten. The hypothesis is that
whatever can be posed this way can be learned the way AlphaGo Zero learned Go: search plays the
games, the outcomes train a network, and the network's first move improves until much of the search
it learned from is no longer needed when the next problem arrives. We call a problem posed this way
an \emph{innovation game} (\cref{sec:formulation}), and we call the claim that a single kind of
learner can play many of them the \emph{innovation-game hypothesis}.

\begin{ezquestion}
Can a model that has watched search solve thousands of innovation games propose a useful solution to
a game it has never seen---in one move, before testing anything? And how many experiments is that
first move worth?
\end{ezquestion}

\paragraph{EurekaZero} We introduce EurekaZero, a proposer that plays innovation games with one
interface (\cref{fig:network,sec:architecture}):
\begin{equation}
\underbrace{z = E_{\phi,d}(s)}_{\code{game\_state\_encoding}}
\qquad\longrightarrow\qquad
\underbrace{h \sim P_\theta(\,\cdot \mid z)}_{\code{solution\_embedding}}
\qquad\longrightarrow\qquad
\underbrace{a = D_{\psi,d}(h)}_{\code{solution}} .
\label{eq:interface}
\end{equation}
A domain adapter encodes the game state $s$, a shared proposer maps it to a solution embedding, and
a typed decoder writes one complete candidate. \emph{Zero} names a protocol, not an absence of
knowledge: the candidate is committed before any candidate-specific execution, scoring or revision,
and only then does an independent verifier evaluate it (\cref{sec:protocol}). Zero does not mean zero
training, zero prior knowledge or zero verification. It means the move is made before the board
answers.

\paragraph{Contributions} We make four contributions, each paired with the evidence that supports it.
\begin{enumerate}
\item \textbf{The innovation game and a commitment protocol} (\cref{sec:formulation,sec:protocol}).
A formulation in which data, histories and specifications are states and programs, decisions and
designs are moves; three output roles; and an auditable boundary---one candidate, no
candidate-specific feedback, no revision---with tracks that separate first proposals from
history-conditioned decisions, sequential control and search.
\item \textbf{EurekaZero} (\cref{sec:architecture,sec:training}). A learned proposer with domain
state adapters, a shared solution embedding and typed decoders with open generative support, trained
by imitation of verified search and then directly on the utility of its first proposal.
\item \textbf{A controlled three-domain evaluation} (\cref{sec:design,sec:results,sec:mechanisms}).
Sealed problem families for learning procedures, next-experiment decisions and engineering parts;
retrieval, catalog, direct-generation, search, learned-baseline and frontier-model controls;
failures retained; complete compute accounting. One proposal improves on the strongest non-search
control in all three domains by \syn[EZ-G01]{\EZdA} to \syn[EZ-G05]{\EZdC} utility and is worth
\syn[EZ-G02]{\EZbeqA} evaluations of search for learning procedures.
\item \textbf{Verified discoveries} (\cref{sec:discoveries}). A novelty ladder from new bytes to
independently replicated mechanisms, and prospective dossiers in which \syn[EZ-G15]{\EZnFive}
proposed mechanisms survive removal and replacement on reserved problems.
\end{enumerate}

\paragraph{What we do not claim} We do not claim that every invention has a machine-readable
specification or a reliable verifier, that EurekaZero discovers new mathematics, that a first
proposal replaces later experiments, or that absence from our archive proves absence from all prior
knowledge. The innovation-game hypothesis is a research program; this paper tests it in three games
(\cref{sec:limits}).

\paragraph{Evidence status} The formulation, protocol, architecture and historical evidence in this
draft are real. The EurekaZero results are synthetic placeholders that show the paper we intend to
earn; \cref{app:gaps} lists the experiment that must replace each one, and \cref{app:history}
reports what the JopeDime program has actually measured so far, including its null results.
